\section{Method: action values, baselines, and conservative allocation}
\label{sec:policy}

\subsection{Outcomes and constraints}

For a policy $d$ and stratum $j\in\{0,1\}$ (no injected fault, injected fault),
define the stratum value and gain over continuation:
\begin{equation}
 V_j(d)=\E[Y_i^{d(Z_i,M_i)}\mid J_i=j],\qquad
 G_j(d)=V_j(d)-V_j(g_0).
 \label{eq:stratum-values}
\end{equation}
Deployment prevalence $p$ mixes the strata:
\begin{equation}
 V(d;p)=pV_1(d)+(1-p)V_0(d),\qquad
 G(d;p)=pG_1(d)+(1-p)G_0(d),
 \label{eq:prevalence-value}
\end{equation}
and cost and invalid-output curves use the same mixture.  The constraints the
first study actually enforces are:
\begin{align}
 V_0(d)&\ge V_0(g_0)-\eta, \label{eq:noninferiority}\\
 C_k(d;p)&\le B_k=\beta_k\,C_k(g_0;p)\quad\text{for each constrained component }k,
 \label{eq:cost-constraint}\\
 R(d;p)&=p\,\E[I_i^{d}\mid J_i=1]+(1-p)\,\E[I_i^{d}\mid J_i=0]\le\xi.
 \label{eq:invalid-constraint}
\end{align}
No-fault noninferiority \eqref{eq:noninferiority} is average no-fault
performance degradation, not an individually harmed fraction; $J=0$ means the
registered fault was not injected and does not guarantee baseline success.
Cost remains a vector: input and output tokens, API expense, CPU and GPU time,
wall time, external calls, storage and network use, human escalation time, and
artifact rewrites are reported; the first study constrains the verifier's
incremental cost components and reports the rest descriptively rather than
inheriting a fixed $K=8$ certification mandate.  The historical
$K=8$/$N=291$ family is a historical registration
(\texttt{protocols/endpoint\_registry.yaml}) and is not inherited.

\subsection{Actionable value}

The actionable value of a checkpoint, defined at \cref{def:actionable} in
\cref{sec:b-theory}, is the largest improvement any available regime achieves
over no intervention.  It cannot detect harm on its own: harm is a property of
an individual regime, and the decision-relevant quantity is whether some
available $g$ carries a negative contrast at a checkpoint.  The two are
reported together.

\subsection{Baselines and the proposed controller}
\label{sec:rl}

All policies share one information set: the pre-treatment features $W=(Z,M)$,
the same permitted model capacity, the same tuning budget and the same branch
data.  The baseline family is:
\begin{enumerate}[label=(\alph*),leftmargin=2em]
\item \textbf{continue everywhere};
\item \textbf{random feasible allocation} at the same capacity;
\item \textbf{always verify where feasible}, reported as a full-budget
  reference and not a matched-count comparator when the feasible rate is below
  one;
\item \textbf{learned failure-risk allocation}: fit $r_S(w)=1-\E[S^0\mid W=w]$
  (or the utility deficit $r_U$ when $U$ is primary) and spend the capacity on
  the highest predicted deficit;
\item \textbf{a transparent predeclared heuristic} over available diagnostics;
\item \textbf{a plain action-conditioned value controller}: fit
  $\hat q(w,g)$ for the two actions and act on
  $\argmax_g \hat q(w,g)$ with feasibility masking --- the closest
  implementation of the nearest prior's action-conditioned selection;
\item \textbf{learned advantage allocation}: fit
  $\hat a(w)=\widehat{\E[O^{g_v}-O^{g_0}\mid W=w]}$ and spend the capacity on
  the largest predicted advantage;
\item \textbf{the same advantage learner with the conservative selection rule
  below}.
\end{enumerate}
Risk and effect learners receive the same permitted $W$, model family, capacity
and tuning budget; the historical $r(Z)$-only score is a secondary baseline, and
richer effect features are never the sole explanation of a claimed target
advantage.  No learner sees condition labels, hidden outcomes, root-identity
shortcuts, sibling results or privileged reference state at deployment.

The learned advantage policy is a one-step masked policy score, not a
regression of a Bellman object:
\begin{equation}
 \widehat J(\theta)=\frac1N\sum_i\sum_{g}\pi_\theta(g\mid Z_i,M_i)\,\widehat A_{ig},
 \qquad
 \pi_\theta(g\mid Z,M)\propto M_g\exp s_\theta(Z,g),
 \label{eq:rl-objective}
\end{equation}
summed over the feasible set only, with checkpoint-relative advantages
$\widehat A_{ig}=\widehat\mu_{ig}-\sum_h\pi_{\mathrm{ref}}(h\mid x_i)\widehat\mu_{ih}$
estimated from branch records.  Full feedback holds conditional on
pre-treatment feasibility, $\Prb(E_{ig}=1\mid A_{ig}=1)=1$, so no
action-assignment positivity must be defended on the feasible set.

\subsection{Conservative selection}
\label{sec:b-conservative}

The first study uses a small deployable policy family, not a cell-wise
certification construction:
\begin{enumerate}[label=(\arabic*),leftmargin=2em]
\item Fit the advantage learner on training data.
\item Construct a small threshold/allocation family on validation data.
\item On independent calibration templates, evaluate paired policy gains,
  no-injected-fault degradation, invalidity and budget.
\item Spend one explicitly allocated simultaneous error budget over the
  family's endpoints, with correctly ranged one-sided bounds.
\item Choose a qualifying policy or fall back to continue everywhere.
\item Report final-test performance of that frozen decision.
\end{enumerate}
A mean-gain guarantee and no-injected-fault noninferiority are population
statements; they do not imply that every checkpoint is helped.  If a
per-context lower-confidence gate is used instead, the conditional mean being
bounded is defined explicitly, coverage across cells/actions/selection is
stated, the comparator is fixed or its selection accounted for, insufficient
support returns to fallback, and no neural confidence score, ensemble spread or
predictive interval is treated as a valid mean-advantage bound without
justification.  A rough conditional improvement statement is permissible only
under its stated assumptions: if simultaneously
$L_g(w)\le Q(w,g)-Q(w,d_0(w))$ and the policy deviates from a feasible
comparator $d_0$ only where $L_g(w)>0$, then the gated policy has nonnegative
population mean advantage on the covered distribution.  This is a standard
pessimistic argument, not a new theorem, and generalization outside the covered
support or to repeated intervention requires new evidence.

Three vocabulary rules are enforced in the paper and in every generated table:
\emph{empirical protection} means observed held-out harm metrics improved;
\emph{certified noninferiority} requires the stated bound, independence, error
allocation and a predeclared tolerance; and \emph{individual harm probability}
is not identified merely from separate stochastic branch marginals.
